
\section{一般光滑性及其推论}

现在，我们能够证实一件你或许早已凭直觉感到的事情。在正\trCharacteristic{}时，这是一条困难的定理，因为它用到了交换代数中一个我们未曾证明的结果。

\textbf{21.6.1.\ \trhTheorem（\trVariety{}的一般光滑性）。}\ ——\ \emph{设 $X$ 是\trPerfectField{} $k$ 上一个\trIntegral{}的\trFiniteType{} $k$-\trScheme，维数为 $n$，那么存在 $X$ 的一个\trDense\trOpenSubset{} $\mathrm{U}$，使得 $\mathrm{U}$（在 $k$ 上）是\trSmooth{}的，且维数为 $n$。}

于是，由\mbox{练习~13.2.R} 结合\mbox{事实~13.8.2}（\trRegular{}的\trLocalization{}是\trRegular{}的），$\mathrm{U}$ 是\trRegular{}的。\mbox{定理~21.6.4} 将把这一结果推广到\trSmooth\trMorphisms，代价是须限制在\trCharacteristic{} $0$ 之中。

若 $k$ 不完全时会出什么问题，可在 \S13.2.9 中找到一个例子。

\emph{证明。}\ $n=0$ 的情形是显然的，故我们设 $n>0$。

我们将证明在\trGenericPoint{}处秩为 $n$。那么，由\trCoherentSheaf{}之秩的\trUpperSemicontinuous{}性（\mbox{练习~14.3.J}），在\trGenericPoint{}的某个\trOpen\trNeighborhood{}内秩必为 $n$，证明便告完成。

于是我们还须验证的是：若 $K$ 是一个维数为 $n$ 的\trIntegral{}的\trFiniteType{} $k$-\trScheme{}的\trFractionField，也就是说（由\mbox{定理~12.2.1}），若 $K/k$ 是一个\trTrDeg{}为 $n$ 的\trFieldExtension，那么 $\Omega_{K/k}$ 是一个 $n$ 维\trVectorSpace。但\trPerfectField{}的每个\trFinitelyGenerated{}扩张都是\trSeparablyGenerated{}的（\mbox{定理~10.5.17}），于是结论由\mbox{练习~21.2.M(b)} 得出。\hfill$\square$

\textbf{21.6.2.\ \mbox{定理~13.8.3} 的证明：\trRegular\trLocalRing{}的\trLocalization{}仍是\trRegular{}的——完全域上\trVariety{}的情形。}\quad \mbox{定理~21.6.1} 使我们得以证明\mbox{定理~13.8.3}（\mbox{事实~13.8.2} 的一个重要情形）。

设 $Y$ 是\trPerfectField{} $k$ 上的一个\trVariety，并在其诸\trClosedPoint{}处\trRegular（于是由\mbox{练习~13.2.R}，$Y$ 是\trSmooth{}的），又设 $\eta$ 是 $Y$ 的一点。我们将证明 $Y$ 在 $\eta$ 处\trRegular。令 $X=\overline{\eta}$。由\mbox{定理~21.6.1}，$X$ 含有一个由\trSmooth{}点组成的\trDense{}（＝非空）\trOpenSubset。把 $X$ 位于该\trOpenSubset{}之外的点丢弃、从而缩小 $Y$ 之后，我们便可假定 $X$ 是\trSmooth{}的。

那么，\mbox{定理~21.2.31}（\trSmooth\trVariety{}的余正规正合列的正合性）蕴涵：$\mathscr{I}/\mathscr{I}^2$ 是一个秩为 $\operatorname{codim}_{X/Y}=\dim\OO_{Y,\eta}$ 的\trLocallyFreeSheaf。

\textbf{21.6.A.（练习）}\quad 设 $\mathfrak{m}$ 是 $\OO_{Y,\eta}$ 的\trMaximalIdeal。把 $\mathscr{I}/\mathscr{I}^2$ 在 $X$ 的\trGenericPoint{} $\eta$ 处的\trStalk{}与 $\mathfrak{m}/\mathfrak{m}^2$ 等同起来。由此完成\mbox{定理~13.8.3} 的证明。\hfill$\square$

\textbf{21.6.B.（棘手练习：\trVariety{}的\trRegular\trLocalRing{}的\trLocalization{}仍是\trRegular{}的——\mbox{定理~13.8.3} 之后所许诺者）}\quad 设 $(A,\mathfrak{m})$ 是一个\trRegular\trLocalRing，同时它是一个\trFinitelyGenerated{} $k$-代数的\trLocalization，其中 $k$ 是完全域。证明：$A$ 在任一\trPrimeIdeal{}处的\trLocalization{}仍是一个\trRegular\trLocalRing。

\textbf{21.6.3.\ \trMorphisms{}的一般光滑性与\trKleiman--\trBertini{}定理。}\quad 现在我们来讨论若干重要的结果，它们都归于「\trMorphisms{}的一般光滑性」这一名目之下。所有这些都要求在\trCharacteristic{} $0$ 中工作。

\textbf{21.6.4.\ \trhTheorem（源上的一般光滑性）。}\ ——\ \emph{(a) 设 $\pi\colon X\to Y$ 是\trIntegral{}的\trNoetherian{}\trScheme{}之间的一个\trDominant\trFiniteType{}\trMorphism，使得 $\operatorname{char}K(Y)$（$=\operatorname{char}K(X)$）为零。那么存在一个非空（＝\trDense）的\trOpenSet{} $\mathrm{U}\subset X$，使得 $\pi|_{\mathrm{U}}$ 是\trSmooth{}的，相对维数为 $\operatorname{trdeg}K(X)/K(Y)$。}

\emph{(b) 设 $k$ 是\trCharacteristic{} $0$ 的域，而 $\pi\colon X\to Y$ 是\trIntegral{}的 $k$-\trVariety{}之间的一个\trDominant\trMorphism。那么存在一个非空（＝\trDense）的\trOpenSet{} $\mathrm{U}\subset X$，使得 $\pi|_{\mathrm{U}}$ 是\trSmooth{}的，相对维数为 $\dim X-\dim Y$。}

\textbf{21.6.C.（练习：\mbox{练习~12.2.D} 中所许诺）}\quad 设 $f_1,\dots,f_n\in k[x_1,\dots,x_n]$，其中 $k$ 是\trCharacteristic{} $0$ 的域。证明：诸 $f_1,\dots,f_n$ \trAlgebraicallyDependent，当且仅当\trJacobianMatrix{}的\trDeterminant{}恒为零。提示：为什么这个练习恰好被安排在\mbox{定理~21.6.4} 之后给出？

\textbf{21.6.5.\ \trhExample。}\quad \mbox{定理~21.6.4} 在正\trCharacteristic{}时失效：考虑\trPurelyInseparable{}扩张 $\mathbb{F}_p(t)/\mathbb{F}_p(t^p)$。同样的问题甚至可能在\trCharacteristic{} $p$ 的\trAlgClosedField{}上出现：考虑 $\mathbb{A}^1_k=\Spec k[t]\to\Spec k[u]=\mathbb{A}^1_k$，由 $u\mapsto t^p$ 给出。%
\footnote{原书此处作「$\mathbb{A}^1_{k'}$」，当为「$\mathbb{A}^1_k$」之误——该态射即 $\Spec k[t]\to\Spec k[u]$，两端同在 $k$ 上——径改。——译者注}

\emph{证明。}\ 令 $n=\operatorname{trdeg}K(X)/K(Y)$。我们可以把 $Y$ 换成一个\trAffine\trOpenSubset{} $\Spec B$，再把 $X$ 换成一个\trAffine\trOpenSubset{} $\Spec A$，于是 $\pi$ 就对应于\trRingMorphism{} $B\to A$。选取 $n$ 个元素 $x_1,\dots,x_n\in A$，使它们构成 $K(A)/K(B)$ 的一组\trTranscendenceBase。（你看出为什么可以选取它们落在 $A$ 中吗？）再选取一些附加的元素 $x_{n+1},\dots,x_N\in A$，使得 $x_1,\dots,x_N$ 把 $A$ 生成为一个 $B$-代数。

对 $i=0,\dots,n$，令 $F_i$ 为 $K(A)$ 在 $K(B)$ 上由 $x_1,\dots,x_i$ 生成的子域。例如，$F_0=K(B)$，$F_N=K(A)$，而当 $i\le n$ 时，$F_i/K(B)$ 是\trPurelyTranscendental\trFieldExtension。对 $i=n+1,\dots,N$，$x_i\in K(A)$ 在 $F_{i-1}$ 上的\trMinimalPolynomial（$F_{i-1}[t]$ 中的一个元素）可以解读为 $\mathfrak{m}_i(x_1,\dots,x_{i-1},t)$ 的形式，其中 $\mathfrak{m}_i\in K(B)[y_1,\dots,y_{i-1},t]$。把 $\mathfrak{m}_i$ 乘以其诸系数分母之乘积，我们便得到 $M_i\in B[y_1,\dots,y_{i-1},t]$，使得 $M_i(x_1,\dots,x_{i-1},t)\in F_{i-1}[t]$ 也是 $x_i\in K(A)$ 在 $F_{i-1}$ 上的\trMinimalPolynomial。于是
\begin{equation*}
K(A)=K(B)(x_1,\dots,x_n)[y_{n+1},\dots,y_N]/I,
\tag{21.6.5.1}\label{eq:generic-smoothness-field}
\end{equation*}
其中
\[
\begin{aligned}
I=\bigl(&M_{n+1}(x_1,\dots,x_n,y_{n+1}),\ M_{n+2}(x_1,\dots,x_n,y_{n+1},y_{n+2}),\\
&\dots,\ M_N(x_1,\dots,x_n,y_{n+1},\dots,y_N)\bigr).
\end{aligned}
\]
定义 $A':=B[y_1,\dots,y_N]/(M_{n+1}(y_1,\dots,y_{n+1}),\dots,M_N(y_1,\dots,y_N))$。

令 $X':=\Spec A'$。设 $Z_1,\dots,Z_s$ 是 $X'$ 的诸\trIrreducibleComponent。（由 $B$ 的诺特性，它们只有有限多个：\mbox{练习~3.6.T} 与\mbox{命题~3.6.15}。）

它们当中恰有一个——比如说 $Z=Z_1$——\trDominant{} $\Spec B$（因为由 \eqref{eq:generic-smoothness-field}，$A'\otimes_B K(B)\cong K(A)$）。令 $\mathrm{U}'\subset X'$ 为\trOpenSubset{} $Z_1\setminus Z_1\cap(Z_2\cup\cdots\cup Z_s)$。

在 $\mathrm{U}'$ 的\trGenericPoint{}处，$M_{n+1}(x_{n+1}),\dots,M_N(x_N)$ 关于 $x_1,\dots,x_N$ 的\trJacobianMatrix{}的\trCorank{}为 $n$。（这里我们用到诸\trMinimalPolynomial{} $M_i$ 的可分性。）于是存在一个\trOpenSubset{} $\mathrm{U}''\subset\mathrm{U}'$，在其上\trJacobianMatrix{}的\trCorank{}为 $n$，从而 $\mathrm{U}''\to\Spec B$ 是\trSmooth{}的，相对维数为 $n$。此外，$\mathrm{U}''$ 是\trIrreducible{}的。现在，$\mathrm{U}''$ 与 $X$ 是\trBirational{}的，故由\mbox{命题~7.5.6}，存在一个\trOpenSubset{} $\mathrm{U}\subset X$，它\trIsomorphic{}于 $\mathrm{U}''$ 的某个\trOpenSubset。这正是我们所要寻找的那个\trOpenSubset。

\textbf{21.6.D.（练习）}\quad 证明\mbox{定理~21.6.4(b)}（利用 (a) 部分）。\hfill$\square$

若 $X$ 还是\trSmooth{}的，情况甚至会更好。

\textbf{21.6.6.\ \trhTheorem（目标上的一般光滑性）。}\ ——\ \emph{设 $\pi\colon X\to Y$ 是 $k$-\trVariety{}之间的一个\trMorphism，其中 $\operatorname{char}k=0$，且 $X$（在 $k$ 上）是\trSmooth{}的。那么存在 $Y$ 的一个\trDense\trOpenSubset{} $\mathrm{U}$，使得 $\pi|_{\pi^{-1}(\mathrm{U})}$ 是一个\trSmooth\trMorphism。}

注意 $\pi^{-1}(\mathrm{U})$ 可能是空的！的确，若 $\pi$ 不是\trDominant{}的，我们就不得不取这样一个 $\mathrm{U}$。

为了证明\mbox{定理~21.6.6}，我们要用一个巧妙的招法。设 $\pi\colon X\to Y$ 是在 $k$ 上\trFiniteType{}的\trScheme{}之间的一个\trMorphism，其中 $\operatorname{char}k=0$。定义
\[
X_r=\bigl\{p\in X\ :\ \operatorname{rank}\bigl(\pi^*\Omega_{Y/k}\to\Omega_{X/k}\bigr)\bigr|_p\le r\bigr\}.
\]

\textbf{21.6.E.（练习）}\quad 证明 $X_r$ 是 $X$ 的一个\trConstructible{}子集。提示：若 $X$ 与 $Y$ 都是\trSmooth{}的，证明秩条件蕴涵 $X_r$ 由「行列式方程」切出，或者使用\trUpperSemicontinuous{}性的某个适当形式。

由\trChevalley{}的\mbox{定理~8.4.2}，$\pi(X_r)$ 是\trConstructible{}的，所以谈论 $\dim\pi(X_r)$ 是有意义的。

\textbf{21.6.7.\ \trhLemma。}\ ——\ \emph{设 $\pi\colon X\to Y$ 是在 $k$ 上\trFiniteType{}的\trScheme{}之间的一个\trMorphism，其中 $\operatorname{char}k=0$，那么 $\dim\pi(X_r)\le r$。}

\emph{证明。}\ 在这个证明中，我们将反复使用「\trClosedPoint{}处的\trZariski{}\trCotangentSpace{}等同于余切层的\trFiber」这一认同，其依据是\mbox{推论~21.2.33}。我们可以把 $X$ 换成 $X_r$ 的某个\trIrreducibleComponent，并把 $Y$ 换成该\trIrreducibleComponent{}的像在 $Y$ 中的\trClosure（带\trReduced\trSubscheme{}结构；见 \S9.4.9）。所得的\trMorphism{}将使整个 $X$ 含于 $X_r$ 之内，所以我们不妨索性设 $X=X_r$。这归结为下面这个线性代数的观察：若线性映射 $\rho\colon V_1\to V_2$ 的秩至多为 $r$，而 $V_1'$ 是 $V_1$ 的一个\trQuotient%
\footnote{原书此处作「a quotient of $V_i$」，当为「a quotient of $V_1$」之误，径改。——译者注}
，且 $\rho$ 把 $V_1'$ 送到 $V_2'$，那么 $\rho$ 在 $V_1'$ 上的\trRestriction{}的秩至多是 $\rho$ 自身的秩。于是我们有一个\trDominant\trMorphism{} $\pi\colon X\to Y$，而我们希望证明 $\dim Y\le r$。由\trVariety{}的一般光滑性（\mbox{定理~21.6.1}），$Y$ 有一个\trDense\trOpenSet{}是\trSmooth{}的，故我们可设 $Y$ 是\trSmooth{}的。由 $\pi\colon X\to Y$ 的源上的一般光滑性，存在一个非空\trOpenSubset{} $\mathrm{U}\subset X$，使得 $\pi\colon\mathrm{U}\to Y$ 是\trSmooth{}的。但对任意一点 $p\in\mathrm{U}$，余切映射 $\pi^*\Omega_{Y|\pi(p)}\to\Omega_{X|p}$ 都是单射（\mbox{练习~21.2.S}），且秩至多为 $r$。取 $p$ 为\trClosedPoint，我们便有 $\dim Y=\dim_{\pi(p)}Y\le\dim\Omega_{Y|\pi(p)}\le r$。\hfill$\square$

要证明这个定理，剩下要做的事就不多了。

\textbf{21.6.8.}\ \emph{\mbox{定理~21.6.6} 的证明。}\quad 先化归到 $Y$ 在 $k$ 上\trSmooth{}的情形（通过限制到更小的\trOpenSet{}来做到，利用 $Y$ 的一般光滑性，即\mbox{定理~21.6.4(b)}）。令 $n=\dim Y$。现在，由\mbox{引理~21.6.7}，$\dim\overline{\pi(X_{n-1})}\le n-1$，故把 $\overline{\pi(X_{n-1})}$ 也从 $Y$ 中去掉。那么，对 $X$ 的每个\trClosedPoint，$\pi^*\Omega_{Y/k}\to\Omega_{X/k}$ 的秩至少为 $n$。但由于 $Y$ 是维数为 $n$ 的\trRegular\trScheme，我们知道：对 $X$ 的每个\trClosedPoint，$\pi^*\Omega_{Y/k}\to\Omega_{X/k}$ 都是单射，从而 $\pi^*\Omega_{Y/k}\to\Omega_{X/k}$ 是\trSheaf{}之间的一个单射映射（你看出其中的缘由吗？）。于是，由\mbox{练习~24.8.G(b)}，$\pi$ 是\trSmooth{}的。\hfill$\square$

\textbf{21.6.9.}\ \extrahard\ \textbf{\trKleiman--\trBertini{}定理。}\quad 对某个「坏\trLocus」的维数作界——这同一想法也可用来证明\mbox{\trKleiman--\trBertini{}定理~21.6.10}（归功于\trKleiman），它在（比如说）计数几何（enumerative geometry）中颇为有用。在这一讨论的全部过程中，$k=\bar{k}$，尽管诸定义与结果都可以推广。设 $G$ 是一个\trGroupVariety（在 $k=\bar{k}$ 上），并且我们有一个 $G$ 在\trVariety{} $X$ 上的作用。若它在诸\trClosedPoint{}上是可迁的，我们就说\textbf{作用是可迁的}。一个更好的定义（你可以证明它与之等价）是：作用\trMorphism{} $G\times X\to X$ 限制在 $X$ 的任一\trClosedPoint{}之上的\trFiber{}上是满的。带有可迁 $G$-作用的\trVariety{} $X$ 称为 $G$ 的一个\textbf{\trHomogeneousSpace}（homogeneous space）。例如，$G$ 通过左\trTranslation{}作用在自身上，而经由这一作用，$G$ 是它自己的一个\trHomogeneousSpace。

\textbf{21.6.F.（容易的练习）}\quad 设 $G$ 是\trCharacteristic{} $0$ 的\trAlgClosedField{} $k$ 上的一个\trGroupVariety。证明：$G$ 的每个\trHomogeneousSpace{} $X$ 都是\trSmooth{}的。（特别地，取 $X=G$，我们看到 $G$ 是\trSmooth{}的。）提示：由\mbox{定理~21.6.4(b)}，$X$ 有一个\trDense\trOpenSet{} $\mathrm{U}$ 是\trSmooth{}的，而 $G$ 在 $X$ 的诸\trClosedPoint{}上可迁地作用，所以我们能用 $\mathrm{U}$ 的诸\trTranslation{}覆盖 $X$。

\textbf{21.6.10.\ \trKleiman--\trBertini{}定理，[Kl2, Thm.~2]。}\ ——\ \emph{设 $X$ 是\trGroupVariety{} $G$（在\trCharacteristic{} $0$ 的域 $k=\bar{k}$ 上）的\trHomogeneousSpace}\footnote{原书此处作「Suppose X is homogeneous space for\ldots」，疑脱不定冠词 a，译文径顺化。——译者注}\emph{，而 $\alpha\colon Y\to X$ 与 $\beta\colon Z\to X$ 是从\trSmooth{} $k$-\trVariety{} $Y$ 与 $Z$ 出发的态射。}

\emph{(a) 那么，存在一个非空\trOpenSubset{} $\mathrm{V}\subset G$，使得对每个 $\sigma\in\mathrm{V}(k)$，由交换图}
\[
\begin{tikzcd}
Y\times_X Z \arrow[r] \arrow[d] & Z \arrow[d, "\beta"] \\
Y \arrow[r, "\sigma\circ\alpha"] & X
\end{tikzcd}
\]
\emph{给出的 $Y\times_X Z$（$Y$ 被 $\sigma$「平移」）是\trSmooth{}的，其维数为 $\dim Y+\dim Z-\dim X$（但可能是空的）。}

\emph{(b) 此外，存在一个非空\trOpenSubset{} $\mathrm{V}\subset G$，使得态射
\begin{equation*}
(G\times_k Y)\times_X Z\longrightarrow G
\tag{21.6.10.1}\label{eq:kleiman-bertini-smooth}
\end{equation*}
是在 $\mathrm{V}$ 上相对维数为 $\dim Y+\dim Z-\dim X$ 的一个\trSmooth\trMorphism。}

第一次听到它时，你应当想到这样的特例：$Y\to X$ 与 $Z\to X$ 都是\trLocallyClosedEmbedding（$Y$ 与 $Z$ 是 $X$ 的\trSmooth\trSubvariety）。在这种情形，\trKleiman--\trBertini{}定理说的是：第二个\trSubvariety{}将与第一个\trSubvariety{}的一个「泛平移」「横截地」（transversely）相交。（不存在负相对维数的\trSmooth\trMorphism{}这种东西。那么，若 $\dim Y+\dim Z-\dim X$ 是负的，会发生什么？）

\emph{证明。}\ 更惬意的做法是「倒着」描述这个证明：设想我们自己会怎样证明它。我们将两次使用一般光滑性。

显然，(b) 蕴涵 (a)，所以我们证明 (b)。

为了证明\trMorphism{} \eqref{eq:kleiman-bertini-smooth} 在 $G$ 的某个非空\trOpenSet{} $\mathrm{V}$ 上是\trSmooth{}的，对 \eqref{eq:kleiman-bertini-smooth} 应用目标上的一般光滑性（\mbox{定理~21.6.6}）便已足够。从而只须证明 $(G\times_k Y)\times_X Z$ 是一个\trSmooth{}的 $k$-\trVariety。现在，$Z$ 在 $k$ 上是\trSmooth{}的，所以只须证明 $(G\times_k Y)\times_X Z\to Z$ 是一个\trSmooth\trMorphism（因为两个\trSmooth\trMorphisms{}的\trComposition{}是\trSmooth{}的，\mbox{练习~13.6.D}）。但这是由 $G\times_k Y\to X$ 经\trBaseChange{}得到的，所以只须证明后一个\trMorphism{}是\trSmooth{}的（因为\trSmoothness{}在\trBaseChange{}下保持）。

现在，$G\times_k Y\to X$ 是一个 $G$-等变（$G$-equivariant）\trMorphism。（所谓「$G$-等变」，我们的意思是两边的 $G$-作用都尊重这个\trMorphism。）由目标上的一般光滑性（\mbox{定理~21.6.6}），它在 $X$ 的某个\trDense\trOpenSubset{}上是\trSmooth{}的。%
\footnote{原书此处作「over a dense open subset X」，当为「over a dense open subset of X」，系脱漏，译文径顺化。——译者注}
但由 $G$-作用在 $X$ 上的可迁性，这个\trMorphism{}处处\trSmooth。

\textbf{21.6.G.（练习）}\quad 改进上面的论证，以证明关于相对维数的那个想要的陈述。

\begin{flushright}
$\square$
\end{flushright}

\textbf{21.6.H.（练习：穷人版的\trKleiman--\trBertini）}\quad 证明\mbox{定理~21.6.10}(a)，但不带关于 $k$ 的假设（\trAlgClosed{}或\trCharacteristic），也不带结论中的\trSmoothness。提示：这是一个关于\trMorphisms{}之\trFiber{}的维数的问题，所以读完 \S12.4 之后你本就能解出它。

\textbf{21.6.I.（练习：改进的\trCharacteristic{} $0$ \trBertini）}\quad 设 $Z$ 是一个\trSmooth{} $k$-\trVariety，其中 $\operatorname{char}k=0$ 且 $k=\bar{k}$。设 $V$ 是 $Z$ 上一个\trFiniteDimensional{}的\trBasePointFree\trLinearSystem，也就是说，由 $Z$ 上某个\trInvertibleSheaf{} $\mathscr{L}$ 的诸\trSection{}组成的一个有限\trVectorSpace。证明：$V$ 的一个泛元素——视为 $Z$ 的一个\trClosedSubscheme——是\trRegular{}的。（更明确地说：每个元素 $s\in V$ 给出 $Z$ 的一个\trClosedSubscheme。那么，对一个泛的 $s$——视为 $\mathbb{P}V$ 中的一点——相应的\trClosedSubscheme{}在 $k$ 上是\trSmooth{}的。）提示：先弄清楚这与\mbox{\trKleiman--\trBertini{}定理~21.6.10} 有什么关系。令 $n=\dim V$，$G=\mathrm{GL}(V)$，$X=\mathbb{P}V^{\vee}$，把\trKleiman--\trBertini{}中的 $Z$ 取作题中的这个 $Z$，而把 $Y$ 取作 $\mathbb{P}V^{\vee}$ 上的「\trUniversalHyperplane」（即\mbox{定义~13.4.1} 的\trIncidenceVariety{} $\mathrm{I}\subset\mathbb{P}V\times\mathbb{P}V^{\vee}$）。

\textbf{21.6.J.（容易的练习）}\quad 把\trCharacteristic{} $0$ 域上的\mbox{\trBertini{}定理~13.4.2} 解读为\mbox{练习~21.6.I} 的一个推论。

在\trCharacteristic{} $0$ 时，\mbox{练习~21.6.I} 是对\trBertini{}定理的一个良好改进。例如，我们不需要 $\mathscr{L}$ 是\trVeryAmple{}的，也不需要 $X$ 是\trProjective{}的。但与\trBertini{}定理不同，\mbox{练习~21.6.I} 在正\trCharacteristic{}时会失效，一维\trLinearSystem{} $\{pQ\ :\ Q\in\mathbb{P}^1\}$ 便是一例。这本质上就是\mbox{例~21.6.5}。（你看出这为什么不与\mbox{\trBertini{}定理~13.4.2} 矛盾吗？）
